\begin{solapartat}
\punts{0,25} Com que el camp magnètic és perpendicular a la superfície de les espires:
\[ \phi = \vec{B}\cdot\vec{S} = BS = B_0 S\cos(\omega t + \varphi_0) \]

\punts{0,5} La força electromotriu en una bobina amb $N$ espires és donada per:
\[ \varepsilon(t) = -N\,\frac{d\phi(t)}{dt} = -N\,S\,\frac{dB(t)}{dt} \]

Si no es té en compte el signe negatiu de la llei de Lenz, cal restar 0,25 p a la
qualificació.

\punts{0,5} $\varepsilon(t) = N\,S\,B_0\,\omega\,\sin(\omega t + \varphi_0)$
\end{solapartat}

\begin{solapartat}
\punts{0,2} $I(t) = \dfrac{\varepsilon(t)}{R} = \dfrac{N\,S\,B_0\,\omega}{R}\,\sin(\omega t + \varphi_0)$

\punts{0,25} $I_{\text{màx}} = \dfrac{N\,S\,B_0\,\omega}{R}$

\punts{0,4} La gràfica correcta és la $A$ ja que una reducció en la freqüència del
moviment implica que l'amplitud disminueixi:
\[ \varepsilon(t) = N\,S\,B_0\,\omega\,\sin(\omega t + \varphi_0) = \varepsilon_{\text{màx}}\sin(\omega t + \varphi_0) \Rightarrow \varepsilon_{\text{màx}} = N\,S\,B_0\,\omega \]
Llavors si $\omega$ disminueix, també ho farà $\varepsilon_{\text{màx}}$.

\punts{0,4} Per altra banda $T = \dfrac{2\pi}{\omega}$, llavors si $\omega$ disminueix el
període augmenta i això es tradueix en un distanciament dels màxims i mínims, les
oscil·lacions s'allarguen en el temps.
\end{solapartat}
